\section{Discussion}

\textbf{Probing is a search problem with a small budget.} A probing agent
learns only what its probes reveal. On Dart, where divergence requires a
specific field to be absent or an integer literal to leave the 64-bit range,
20 probes chosen by the model's prior found a divergence on one seed in five.
Table~\ref{tab:coverage} classifies all 300 probes by what they change
relative to a valid record. The agents probed the same way on both targets:
mostly valid documents, then one wrong type, an unknown enum constant or an
extra key. On neither target did any probe remove a required top-level field,
use an \texttt{id} near a 53- or 64-bit boundary, or send text that is not
JSON, although probes were raw text precisely to allow that. Wrong types,
unknown enum constants and extra keys never diverged on Dart; on Kotlin they
are exactly where the libraries disagree. The difference between the targets is the density of
divergence near the inputs the model finds natural, not the agent's
behaviour.

\begin{table}[t]
\centering
\caption{What the probes changed, relative to a valid record (a probe can
be in several rows). Seeds: on how many seeds at least one probe did it.}
\label{tab:coverage}
\footnotesize
\begin{tabular}{@{}lrrrr@{}}
\toprule
& \multicolumn{2}{c}{Dart (100 probes)} & \multicolumn{2}{c}{Kotlin (200)} \\
Change & probes & seeds & probes & seeds \\
\midrule
nothing at top level     & 68 & 5/5 & 117 & 10/10 \\
wrong type               & 19 & 5/5 & 48 & 10/10 \\
unknown enum constant    & 6  & 5/5 & 11 & 10/10 \\
extra key                & 6  & 5/5 & 14 & 9/10 \\
null for non-null field  & 1  & 1/5 & 10 & 8/10 \\
nested \texttt{child} changed & 19 & 5/5 & 42 & 10/10 \\
missing required field   & 0  & 0/5 & 0  & 0/10 \\
\texttt{id} near 53/64-bit boundary & 0 & 0/5 & 0 & 0/10 \\
not JSON                 & 0  & 0/5 & 0  & 0/10 \\
\bottomrule
\end{tabular}
\end{table}
The agent's ability to characterize a target is therefore bounded by the
same thing that bounds the fuzzer: the density of interesting behaviour near
the inputs the model finds natural. This is the conditional result in our
title question. Self-acquired knowledge helps when an unguided search would
have found some divergence anyway, and gives nothing where knowledge is most
needed.

\textbf{Negative evidence is misleading.} A probe summary that observed no
divergence does not say ``I did not look at missing fields''; it says that
the implementations behave identically, sometimes with an overgeneralization
that is false. The loop then receives confident but empty knowledge. A
practical probing agent should report what it did \emph{not} test, or be
required to cover categories of perturbation, such as each field missing,
each field null and each numeric boundary, before summarizing. Such a checklist
is itself domain knowledge, so it moves the question rather than removing
it, but it is generic across schemas and cheap to write once.

\textbf{Code reading is reliable but incomplete.} The code arm never failed,
because the facts are in the code whether or not the agent thinks to test
them. It stated the \texttt{num.toInt()} difference in every summary and
found \texttt{RAAR} in every run. It missed \texttt{built\_value}'s
empty-list default, which lives in generated builder code rather than in
the decoding call, and it never connected \texttt{toInt()} to
\texttt{jsonDecode}'s treatment of large literals, which lives in the SDK and
was not in its reading list. Reading more code would help, at the cost of the
fixed budget. On Kotlin the arm could not be run at all: the equivalent logic
is reflective or compiler-generated.

\textbf{Combining sources.} Code reading and probing fail differently. Code
reading states what the implementation does wherever the logic is visible,
but misses behaviour that lives outside its reading list, such as the SDK's
number parsing. Probing observes behaviour end to end, including such
interactions, but only where it looks. A natural combination, which we did
not evaluate, is to let code reading propose where to probe, and probing
confirm or refute what code reading claims. That would also address the
overgeneralization we observed, because a claim of agreement would then have
to name the probes that support it.

\textbf{For practitioners.} Where a human characterization exists, use it:
it was the best source by a wide margin and costs nothing at run time. Where
source is available, a code-reading summary gives a dependable improvement
for a few cents. A probing agent is worth running first to learn whether
disagreement is common; if its probes find none, treat its summary as
uninformative, not as evidence of agreement.

\subsection{Threats to validity}

\textbf{Sample size.} Each arm has five seeds, the smallest size at which an
exact Mann--Whitney test can reach $p < 0.05$. The primary Dart result is a
non-detection, and the bimodal probe distribution means a larger sample
would most likely estimate how often probing hits, not remove the
variance. The Kotlin result relies on a pre-declared replication and is
weaker on RFC-valid documents alone ($p \approx 0.075$).

\textbf{Knowledge quality is uncontrolled.} Summaries are capped at the
same length but differ in accuracy, which is the variable under study. The
cap was applied by cutting in 19 of 20 acquisitions, which drops the last
observations; a different enforcement rule could change individual
summaries.

\textbf{One model family and one probe budget.} All agents and the loop use
\texttt{gpt-4.1-mini}, and the probe budget of 20 is a choice. A stronger
model or a larger budget might probe more widely; the previous study found
that a larger model alone did not close the gap.

\textbf{Post-data change.} The Kotlin harness's validity label was found to
be lenient after the data existed. It affects only a secondary metric, which
we recompute with a strict parser; the harness is unchanged so that the
committed runs stay reproducible. The output-token limit was raised before
any experiment data, for every arm.

\textbf{Non-determinism and model drift.} The API is not seedable, so a
seed labels a run rather than fixing it; re-running reproduces the procedure,
not the documents. All prompts, responses, generators and per-document
results are committed so that every reported number can be recomputed without
calling the API. A hosted model may also change over time; all runs were
made within two days, and the model identifier is recorded in each
manifest.

\textbf{Answer keys and generality.} Pattern recall is measured against two
known Dart patterns and one known Kotlin behaviour; a source could help in
ways these keys do not capture. Both targets decode one small record schema
from JSON, with four paths each. The conditional result, that probing helps
where divergence is dense, is a hypothesis for other formats and schemas, not
a finding about them.

\textbf{Harness fidelity.} The Dart arms reuse the previous study's harness
binary and loop; a test rebuilds all 100 of its committed prompts byte for
byte, and its committed results summarize identically under the ported
code. The Kotlin harness is new; each behaviour the paper reports was
reproduced with a minimal hand-written document, and the control-character
behaviour also in a standalone program outside the harness, on two library
versions.

\textbf{Construct.} ``Divergence'' counts documents, and on Kotlin many of
them express the same few documented behaviours. Pattern recall and the
triage are reported alongside the rate for this reason.
